\documentclass{amsart}
\usepackage[UTF8,fontset=windows]{ctex}
\usepackage[a4paper,margin=24mm]{geometry}
% For dealing with references we use the comment environment
\usepackage{verbatim}
\newenvironment{reference}{\comment}{\endcomment}
%\newenvironment{reference}{}{}
\newenvironment{slogan}{\comment}{\endcomment}
\newenvironment{history}{\comment}{\endcomment}

% For commutative diagrams we use Xy-pic
\usepackage[all]{xy}

% We use 2cell for 2-commutative diagrams.
\xyoption{2cell}
\UseAllTwocells

% We use multicol for the list of chapters between chapters
\usepackage{multicol}

% This is generally recommended for better output
\usepackage{lmodern}
\usepackage[T1]{fontenc}

% For cross-file-references

% Package for hypertext links:
\usepackage{hyperref}

% For any local file, say "hello.tex" you want to link to please

% Theorem environments.
%
\theoremstyle{plain}
\newtheorem{theorem}[subsection]{Theorem}
\newtheorem{proposition}[subsection]{Proposition}
\newtheorem{lemma}[subsection]{Lemma}

\theoremstyle{definition}
\newtheorem{definition}[subsection]{Definition}
\newtheorem{example}[subsection]{Example}
\newtheorem{exercise}[subsection]{Exercise}
\newtheorem{situation}[subsection]{Situation}

\theoremstyle{remark}
\newtheorem{remark}[subsection]{Remark}
\newtheorem{remarks}[subsection]{Remarks}

\numberwithin{equation}{subsection}

% Macros
%
\def\lim{\mathop{\mathrm{lim}}\nolimits}
\def\colim{\mathop{\mathrm{colim}}\nolimits}
\def\Spec{\mathop{\mathrm{Spec}}}
\def\Hom{\mathop{\mathrm{Hom}}\nolimits}
\def\Ext{\mathop{\mathrm{Ext}}\nolimits}
\def\SheafHom{\mathop{\mathcal{H}\!\mathit{om}}\nolimits}
\def\SheafExt{\mathop{\mathcal{E}\!\mathit{xt}}\nolimits}
\def\Sch{\mathit{Sch}}
\def\Mor{\mathop{\mathrm{Mor}}\nolimits}
\def\Ob{\mathop{\mathrm{Ob}}\nolimits}
\def\Sh{\mathop{\mathit{Sh}}\nolimits}
\def\NL{\mathop{N\!L}\nolimits}
\def\CH{\mathop{\mathrm{CH}}\nolimits}
\def\proetale{{pro\text{-}\acute{e}tale}}
\def\etale{{\acute{e}tale}}
\def\QCoh{\mathit{QCoh}}
\def\Ker{\mathop{\mathrm{Ker}}}
\def\Im{\mathop{\mathrm{Im}}}
\def\Coker{\mathop{\mathrm{Coker}}}
\def\Coim{\mathop{\mathrm{Coim}}}

% Boxtimes
%
\DeclareMathSymbol{\boxtimes}{\mathbin}{AMSa}{"02}

%
% Macros for moduli stacks/spaces
%
\def\QCohstack{\mathcal{QC}\!\mathit{oh}}
\def\Cohstack{\mathcal{C}\!\mathit{oh}}
\def\Spacesstack{\mathcal{S}\!\mathit{paces}}
\def\Quotfunctor{\mathrm{Quot}}
\def\Hilbfunctor{\mathrm{Hilb}}
\def\Curvesstack{\mathcal{C}\!\mathit{urves}}
\def\Polarizedstack{\mathcal{P}\!\mathit{olarized}}
\def\Complexesstack{\mathcal{C}\!\mathit{omplexes}}
% \Pic is the operator that assigns to X its picard group, usage \Pic(X)
% \Picardstack_{X/B} denotes the Picard stack of X over B
% \Picardfunctor_{X/B} denotes the Picard functor of X over B
\def\Pic{\mathop{\mathrm{Pic}}\nolimits}
\def\Picardstack{\mathcal{P}\!\mathit{ic}}
\def\Picardfunctor{\mathrm{Pic}}
\def\Deformationcategory{\mathcal{D}\!\mathit{ef}}

\hypersetup{hidelinks}
\widowpenalty=10000\clubpenalty=10000\displaywidowpenalty=10000
\begin{document}
\title[Stacks Project, Tag 09F1]{等特征 p 的弱消分歧构造}
\maketitle
\noindent\textbf{Counting scope.} The primary problem is Tag 09F1, marked below. All preceding numbered statements are supporting author context and are not counted as additional problems. This is a local pilot with a temporary ID, not production admission.


\noindent Copyright (C) 2005--2025 Johan de Jong. Stacks Project authors. Source: \href{https://stacks.math.columbia.edu/tag/09F1}{Tag 09F1}.

\noindent Difficulty H2: 本题为满足残余域交条件的DVR扩张构造全分歧有限Galois扩张，使度p扩张得到弱解。证明展开Artin–Schreier参数，消去互为p幂倍的极点指数，再选择全分歧扩张使唯一主导极点落入指定分歧分类。预备分歧分类和Cohen截面可引用，但参数修正、归纳及扩张选择仍须本证明完成，判H2。.

\noindent History: excerpt from revision local cached snapshot; exact chapter hashes; revision unknown, more-algebra.tex, lines 34817--34974. Only reference presentation was adapted. Licensed under GNU FDL 1.2 or later; no invariant sections or cover texts. See accompanying COPYING.

\begin{definition}
\label{definition-solution}
Let $A \to B$ be an extension of discrete valuation rings with fraction
fields $K \subset L$.
\begin{enumerate}
\item We say a finite field extension $K_1/K$ is a
{\it weak solution for $A \subset B$} if all the extensions
$(A_1)_{\mathfrak m_i} \subset (B_1)_{\mathfrak m_{ij}}$ of
Remark \ref{remark-construction} are weakly unramified.
\item We say a finite field extension $K_1/K$ is a
{\it solution for $A \subset B$} if each extension
$(A_1)_{\mathfrak m_i} \subset (B_1)_{\mathfrak m_{ij}}$ of
Remark \ref{remark-construction} is formally smooth in
the $\mathfrak m_{ij}$-adic topology.
\end{enumerate}
We say a solution $K_1/K$ is a {\it separable solution}
if $K_1/K$ is separable.
\end{definition}

\begin{remark}
\label{remark-construction}
Let $A \to B$ be an extension of discrete valuation rings with fraction
fields $K \subset L$. Let $K_1/K$ be a finite extension of
fields. Let $A_1 \subset K_1$ be the integral closure of $A$ in $K_1$.
On the other hand, let $L_1 = (L \otimes_K K_1)_{red}$. Then $L_1$ is a
nonempty finite product of finite field extensions of $L$. Let $B_1$ be
the integral closure of $B$ in $L_1$. We obtain compatible commutative
diagrams
$$
\vcenter{
\xymatrix{
L \ar[r] & L_1 \\
K \ar[u] \ar[r] & K_1 \ar[u]
}
}
\quad\text{and}\quad
\vcenter{
\xymatrix{
B \ar[r] & B_1 \\
A \ar[u] \ar[r] & A_1 \ar[u]
}
}
$$
In this situation we have the following
\begin{enumerate}
\item By Algebra, Lemma \href{https://stacks.math.columbia.edu/tag/09IG}{lemma integral closure Dedekind (Tag 09IG)}
the ring $A_1$ is a Dedekind domain and $B_1$ is a finite product of
Dedekind domains.
\item Note that $L \otimes_K K_1 = (B \otimes_A A_1)_\pi$ where $\pi \in A$
is a uniformizer and that $\pi$ is a nonzerodivisor on $B \otimes_A A_1$. 
Thus the ring map $B \otimes_A A_1 \to B_1$ is integral with kernel
consisting of nilpotent elements. Hence $\Spec(B_1) \to \Spec(B \otimes_A A_1)$
is surjective on spectra
(Algebra, Lemma \href{https://stacks.math.columbia.edu/tag/00GQ}{lemma integral overring surjective (Tag 00GQ)}).
The map $\Spec(B \otimes_A A_1) \to \Spec(A_1)$ is surjective as
$A_1/\mathfrak m_A A_1 \to
B/\mathfrak m_AB \otimes_{\kappa_A} A_1/\mathfrak m_A A_1$
is an injective ring map with $A_1/\mathfrak m_A A_1$ Artinian.
We conclude that $\Spec(B_1) \to \Spec(A_1)$ is surjective.
\item Let $\mathfrak m_i$, $i = 1, \ldots n$ with $n \geq 1$ be the
maximal ideals of $A_1$. For each $i = 1, \ldots, n$ let
$\mathfrak m_{ij}$, $j = 1, \ldots, m_i$ with $m_i \geq 1$
be the maximal ideals of $B_1$ lying over $\mathfrak m_i$. We obtain diagrams
$$
\xymatrix{
B \ar[r] & (B_1)_{\mathfrak m_{ij}} \\
A \ar[u] \ar[r] & (A_1)_{\mathfrak m_i} \ar[u]
}
$$
of extensions of discrete valuation rings.
\item If $A$ is henselian (for example complete), then $A_1$ is a
discrete valuation ring, i.e., $n = 1$.
Namely, $A_1$ is a union of finite extensions of $A$ which are domains,
hence local by Algebra, Lemma \href{https://stacks.math.columbia.edu/tag/04GH}{lemma finite over henselian (Tag 04GH)}.
\item If $B$ is henselian (for example complete), then $B_1$
is a product of discrete valuation rings, i.e., $m_i = 1$ for
$i = 1, \ldots, n$.
\item If $K \subset K_1$ is purely inseparable, then $A_1$ and $B_1$
are both discrete valuation rings, i.e., $n = 1$ and $m_1 = 1$.
This is true because for every $b \in B_1$ a $p$-power power of $b$
is in $B$, hence $B_1$ can only have one maximal ideal.
\item If $K \subset K_1$ is finite separable, then $L_1 = L \otimes_K K_1$
and is a finite product of finite separable extensions too. Hence
$A \subset A_1$ and $B \subset B_1$ are finite by
Algebra, Lemma
\href{https://stacks.math.columbia.edu/tag/032L}{lemma Noetherian normal domain finite separable extension (Tag 032L)}.
\item If $A$ is Nagata, then $A \subset A_1$ is finite.
\item If $B$ is Nagata, then $B \subset B_1$ is finite.
\end{enumerate}
\end{remark}

\begin{lemma}
\label{lemma-pre-characteristic-p-case}
Let $A$ be a discrete valuation ring with fraction field $K$ of characteristic
$p > 0$. Let $\xi \in K$. Let $L$ be an extension of $K$ obtained by
adjoining a root of $z^p - z = \xi$. Then $L/K$ is Galois and one of the
following happens
\begin{enumerate}
\item $L = K$,
\item $L/K$ is unramified with respect to $A$ of degree $p$,
\item $L/K$ is totally ramified with respect to $A$
with ramification index $p$, and
\item the integral closure $B$ of $A$ in $L$ is a discrete valuation ring,
$A \subset B$ is weakly unramified, and $A \to B$ induces a purely inseparable
residue field extension of degree $p$.
\end{enumerate}
Let $\pi$ be a uniformizer of $A$. We have the following implications:
\begin{enumerate}
\item[(A)] If $\xi \in A$, then we are in case (1) or (2).
\item[(B)] If $\xi = \pi^{-n}a$ where $n > 0$ is not divisible by
$p$ and $a$ is a unit in $A$, then we are in case (3)
\item[(C)] If $\xi = \pi^{-n} a$ where $n > 0$ is divisible by $p$ and
the image of $a$ in $\kappa_A$ is not a $p$th power, then we are in case (4).
\end{enumerate}
\end{lemma}

\begin{proof}
The extension is Galois of order dividing $p$ by the discussion in
Fields, Section \href{https://stacks.math.columbia.edu/tag/09I7}{section Artin Schreier (Tag 09I7)}.
It immediately follows from the discussion in
Section \href{https://stacks.math.columbia.edu/tag/09E3}{section ramification (Tag 09E3)} that we are in one of the cases (1) -- (4)
listed in the lemma.

\medskip\noindent
Case (A). Here we see that $A \to A[x]/(x^p - x - \xi)$ is a finite
\'etale ring extension. Hence we are in cases (1) or (2).

\medskip\noindent
Case (B). Write $\xi = \pi^{-n}a$ where $p$ does not divide $n$.
Let $B \subset L$ be the integral closure of $A$ in $L$.
If $C = B_\mathfrak m$ for some maximal ideal $\mathfrak m$,
then it is clear that $p \text{ord}_C(z) = -n \text{ord}_C(\pi)$.
In particular $A \subset C$ has ramification index divisible by $p$.
It follows that it is $p$ and that $B = C$.

\medskip\noindent
Case (C). Set $k = n/p$. Then we can rewrite the equation as
$$
(\pi^kz)^p - \pi^{n - k} (\pi^kz) = a
$$
Since $A[y]/(y^p - \pi^{n - k}y - a)$ is a discrete valuation ring
weakly unramified over $A$, the lemma follows.
\end{proof}


\begin{lemma}
\label{lemma-cohen}
Let $A$ be a local ring annihilated by a prime $p$ whose maximal ideal is
nilpotent. There exists a ring map $\sigma : \kappa_A \to A$
which is a section to the residue map $A \to \kappa_A$. If $A \to A'$ is
a local homomorphism of local rings, then we can choose a similar
ring map $\sigma' : \kappa_{A'} \to A'$ compatible with $\sigma$ provided
that the extension $\kappa_{A'}/\kappa_A$ is separable.
\end{lemma}

\begin{lemma}
\label{lemma-make-degree-q-extension}
Let $A$ be a discrete valuation ring with uniformizer $\pi$. If the residue
characteristic of $A$ is $p > 0$, then for every $n > 1$ and $p$-power $q$
there exists a degree $q$ separable extension $L/K$
totally ramified with respect to $A$
such that the integral closure $B$ of $A$ in $L$ has ramification index
$q$ and a uniformizer $\pi_B$ such that
$\pi_B^q = \pi + \pi^n b$ and $\pi_B^q = \pi + (\pi_B)^{nq}b'$
for some $b, b' \in B$.
\end{lemma}

\section*{Primary problem: Tag 09F1}
\begin{lemma}
\label{lemma-characteristic-p-case}
Let $A \subset B \subset C$ be extensions of discrete valuation rings
with fractions fields $K \subset L \subset M$. Assume
\begin{enumerate}
\item $A \subset B$ weakly unramified,
\item the characteristic of $K$ is $p$,
\item $M$ is a degree $p$ Galois extension of $L$, and
\item $\kappa_A = \bigcap_{n \geq 1} \kappa_B^{p^n}$.
\end{enumerate}
Then there exists a finite Galois extension $K_1/K$
totally ramified with respect to $A$
which is a weak solution for $A \to C$.
\end{lemma}

\begin{proof}
Since the characteristic of $L$ is $p$ we know that $M$ is an Artin-Schreier
extension of $L$ (Fields, Lemma \href{https://stacks.math.columbia.edu/tag/09DY}{lemma Artin Schreier (Tag 09DY)}).
Thus we may pick $z \in M$, $z \not \in L$ such that
$\xi = z^p - z \in L$. Choose $n \geq 0$ such that $\pi^n\xi \in B$.
We pick $z$ such that $n$ is minimal. If $n = 0$, then $M/L$ is unramified
with respect to $B$ (Lemma \ref{lemma-pre-characteristic-p-case}) and
we are done. Thus we have $n > 0$.

\medskip\noindent
Assumption (4) implies that $\kappa_A$ is perfect. Thus we may
choose compatible ring maps $\overline{\sigma} : \kappa_A \to A/\pi^n A$ and
$\overline{\sigma} : \kappa_B \to B/\pi^n B$ as in
Lemma \ref{lemma-cohen}. We lift the second of these to a
map of sets $\sigma : \kappa_B \to B$\footnote{If $B$ is complete, then
we can choose $\sigma$ to be a ring map. If $A$ is also complete and
$\sigma$ is a ring map, then $\sigma$ maps $\kappa_A$ into $A$.}.
Then we can write
$$
\xi = \sum\nolimits_{i = n, \ldots, 1} \sigma(\lambda_i) \pi^{-i} + b
$$
for some $\lambda_i \in \kappa_B$ and $b \in B$. Let
$$
I = \{i \in \{n, \ldots, 1\} \mid \lambda_i \in \kappa_A\}
$$
and
$$
J = \{j \in \{n, \ldots, 1\} \mid \lambda_i \not \in \kappa_A\}
$$
We will argue by induction on the size of the finite set $J$.

\medskip\noindent
The case $J = \emptyset$. Here for all $i \in \{n, \ldots, 1\}$ we have
$\sigma(\lambda_i) = a_i + \pi^n b_i$ for some $a_i \in A$ and $b_i \in B$
by our choice of $\sigma$. Thus
$\xi = \pi^{-n} a + b$ for some $a \in A$ and $b \in B$.
If $p | n$, then we write $a = a_0^p + \pi a_1$ for some $a_0, a_1 \in A$
(as the residue field of $A$ is perfect). We compute
$$
(z - \pi^{-n/p}a_0)^p - (z - \pi^{-n/p}a_0) =
\pi^{-(n - 1)}(a_1 + \pi^{n - 1 - n/p}a_0) + b'
$$
for some $b' \in B$. This would contradict the minimality of $n$. Thus $p$
does not divide $n$. Consider the degree $p$ extension $K_1$ of $K$ given
by $w^p - w = \pi^{-n}a$. By Lemma \ref{lemma-pre-characteristic-p-case}
this extension is Galois and totally ramified with respect to $A$.
Thus $L_1 = L \otimes_K K_1$ is a field and $A_1 \subset B_1$
is weakly unramified
(Lemma \href{https://stacks.math.columbia.edu/tag/09ER}{lemma weakly unramified goes up along totally ramified (Tag 09ER)}).
By Lemma \ref{lemma-pre-characteristic-p-case}
the ring $M_1 = M \otimes_K K_1$ is either a product of $p$ copies
of $L_1$ (in which case we are done) or a field extension of $L_1$
of degree $p$. Moreover, in the second case, either $C_1$ is weakly unramified
over $B_1$ (in which case we are done) or $M_1/L_1$ is degree $p$,
Galois, and totally ramified with respect to $B_1$.
In this last case the extension $M_1/L_1$
is generated by the element $z - w$ and
$$
(z - w)^p - (z - w) = z^p - z - (w^p - w) = b
$$
with $b \in B$ (see above). Thus by Lemma \ref{lemma-pre-characteristic-p-case}
once more the extension $M_1/L_1$ is unramified with respect to $B_1$
and we conclude that $K_1$ is a weak solution for $A \to C$.
From now on we assume $J \not = \emptyset$.

\medskip\noindent
Suppose that $j', j \in J$ such that $j' = p^r j$ for some
$r > 0$. Then we change our choice of $z$ into
$$
z' = z -
(\sigma(\lambda_j) \pi^{-j} + \sigma(\lambda_j^p) \pi^{-pj} + \ldots +
\sigma(\lambda_j^{p^{r - 1}}) \pi^{-p^{r - 1}j})
$$
Then $\xi$ changes into $\xi' = (z')^p - (z')$ as follows
$$
\xi' =
\xi - \sigma(\lambda_j) \pi^{-j} + \sigma(\lambda_j^{p^r}) \pi^{-j'}
+ \text{something in }B
$$
Writing
$\xi' = \sum\nolimits_{i = n, \ldots, 1} \sigma(\lambda'_i) \pi^{-i} + b'$
as before we find that
$\lambda'_i = \lambda_i$ for $i \not = j, j'$ and $\lambda'_j = 0$.
Thus the set $J$ has gotten smaller.
By induction on the size of $J$ we may assume no such pair $j, j'$ exists.
(Please observe that in this procedure we may get thrown back into the case
that $J = \emptyset$ we treated above.)

\medskip\noindent
For $j \in J$ write $\lambda_j = \mu_j^{p^{r_j}}$ for some $r_j \geq 0$ and
$\mu_j \in \kappa_B$ which is not a $p$th power. This is possible by our
assumption (4). Let $j \in J$ be the unique index such that $j p^{-r_j}$
is maximal. (The index is unique by the result of the preceding paragraph.)
Choose $r > \max(r_j + 1)$ and such that $j p^{r - r_j} > n$ for $j \in J$.
Choose a separable extension $K_1/K$ totally ramified with respect to $A$
of degree $p^r$ such that the corresponding discrete valuation ring
$A_1 \subset K_1$ has uniformizer $\pi'$ with
$(\pi')^{p^r} = \pi + \pi^{n + 1}a$ for some $a \in A_1$
(Lemma \ref{lemma-make-degree-q-extension}).
Observe that $L_1 = L \otimes_K K_1$ is a field and that
$L_1/L$ is totally ramified with respect to $B$
(Lemma \href{https://stacks.math.columbia.edu/tag/09ER}{lemma weakly unramified goes up along totally ramified (Tag 09ER)}).
Computing in the integral closure $B_1$ we get
$$
\xi = \sum\nolimits_{i \in I} \sigma(\lambda_i) (\pi')^{-i p^r} +
\sum\nolimits_{j \in J} \sigma(\mu_j)^{p^{r_j}} (\pi')^{-j p^r} + b_1
$$
for some $b_1 \in B_1$. Note that $\sigma(\lambda_i)$ for $i \in I$
is a $q$th power modulo $\pi^n$, i.e., modulo $(\pi')^{n p^r}$.
Hence we can rewrite the above as
$$
\xi = \sum\nolimits_{i \in I} x_i^{p^r} (\pi')^{-i p^r} +
\sum\nolimits_{j \in J} \sigma(\mu_j)^{p^{r_j}} (\pi')^{- j p^r}
+ b_1
$$
As in the previous paragraph we change our choice of $z$ into
\begin{align*}
z' & = z \\
& -
\sum\nolimits_{i \in I}
\left(x_i (\pi')^{-i} + \ldots + x_i^{p^{r - 1}} (\pi')^{-i p^{r - 1}}\right)
\\
& -
\sum\nolimits_{j \in J}
\left(
\sigma(\mu_j) (\pi')^{- j p^{r - r_j}}
+ \ldots +
\sigma(\mu_j)^{p^{r_j - 1}} (\pi')^{- j p^{r - 1}}
\right)
\end{align*}
to obtain
$$
(z')^p - z' =
\sum\nolimits_{i \in I} x_i (\pi')^{-i} +
\sum\nolimits_{j \in J} \sigma(\mu_j) (\pi')^{- j p^{r - r_j}} + b_1'
$$
for some $b'_1 \in B_1$.
Since there is a unique $j$ such that $j p^{r - r_j}$ is maximal
and since $j p^{r - r_j}$ is bigger than $i \in I$ and divisible
by $p$, we see that $M_1 / L_1$ falls into case (C) of
Lemma \ref{lemma-pre-characteristic-p-case}.
This finishes the proof.
\end{proof}
\end{document}
